% Volume VIII: Computational Neuroscience Mathematics
% Continuity note for Chapter 50.
% Previous dependencies: Chapter 15's Bayes theorem, conditional independence, posterior risk, and Bayes actions; Chapter 18's KL divergence, cross-entropy, and information; Chapter 22's variational objectives; Chapter 27's stochastic-process and filtering notation when available; Chapter 30's dynamic programming when available; Chapter 47's Bayesian decoding; Chapter 49's recurrent dynamics and free-energy vocabulary.
% New durable concepts: evidence accumulation; drift-diffusion models; log-likelihood ratio accumulation; speed-accuracy tradeoff; Bayesian cue combination; sequential Bayes; generative models as latent cause to observation; predictive coding as precision-weighted prediction-error minimization; active inference as policy selection under model-based prediction and preferences; limits of Bayesian-brain analogies.
% Forward hooks: Chapter 51 uses these models in neural data analysis and model comparison; Chapter 52 discusses bridges and limits between biological and artificial intelligence; Chapter 53 integrates Bayes and predictive coding as a capstone case study; Chapter 54 summarizes Bayesian, SDE, and variational notation.
% Notation commitments: Hidden causes or states are Z or X_t; observations are Y or O_t; actions are A_t; policies are \pi; beliefs are probability distributions p(z\mid y) or q(z); prediction errors are \varepsilon; precisions are inverse variances or inverse covariance matrices.
\section{Chapter 50: Decision-Making, Bayesian Brain, and Predictive Coding}
\label{ch:decision-making-bayesian-brain-predictive-coding}

The previous chapters developed encoding, decoding, population dynamics, and recurrent circuits. Chapter 50 asks how perception and action can be modeled as inference under uncertainty. A decision-maker rarely observes the relevant state directly. It accumulates noisy evidence, combines prior knowledge with sensory measurements, updates beliefs over time, and chooses actions that shape future observations.

The phrase ``Bayesian brain'' should be handled carefully. Bayes' theorem is an exact identity about probability models. Predictive coding and active inference are model classes and theoretical proposals. They can generate useful equations, experiments, and AI algorithms, but they do not prove that every neural circuit literally computes the exact posterior of a specified model. This chapter states the mathematics first and separates it from biological interpretation.

\paragraph{Chapter problem.}
How can evidence accumulation, Bayesian perception, predictive coding, and active inference be written as precise models of decision, update, prediction error, and action?

\paragraph{Section map.}
Section 50.1 introduces drift-diffusion and log-likelihood evidence accumulation. Section 50.2 derives Bayesian perception, cue combination, and sequential Bayes. Section 50.3 derives predictive-coding updates from a simple linear-Gaussian generative model. Section 50.4 frames active inference and control as policy-dependent prediction, preference, and information-seeking under a generative model.

\subsection{50.1 Drift-diffusion and evidence accumulation}

\paragraph{Motivating question.}
How can a noisy stream of evidence be accumulated until a decision threshold is reached?

\paragraph{Beginner setup.}
Many decisions unfold over time. A subject may view a noisy motion stimulus and decide left or right. Each moment supplies a little evidence. The decision variable accumulates that evidence until it reaches a boundary. The boundary crossing determines both the choice and the reaction time.

The drift-diffusion model, or DDM, is the simplest continuous-time model of this idea:
\[
dX_t=\mu\,dt+\sigma\,dW_t.
\]
Here \(X_t\) is accumulated evidence, \(\mu\) is drift, \(\sigma\) is noise scale, and \(W_t\) is Brownian motion. Upper and lower boundaries represent the two choices.

\paragraph{Formal setup.}
Let \(X_t\in\mathbb R\) be a decision variable. The drift-diffusion model is the stochastic differential equation
\[
\boxed{
dX_t=\mu\,dt+\sigma\,dW_t,\qquad X_0=x_0,
}
\]
with absorbing boundaries at \(+a\) and \(-a\), where \(a>0\). The decision time is
\[
T=\inf\{t\geq0:X_t\in\{-a,+a\}\}.
\]
The choice is \(+1\) if \(X_T=+a\) and \(-1\) if \(X_T=-a\).

In discrete time, with time step \(\Delta t\), an Euler approximation is
\[
X_{k+1}=X_k+\mu\Delta t+\sigma\sqrt{\Delta t}\,\xi_k,
\qquad
\xi_k\sim\mathcal N(0,1).
\]
The \(\sqrt{\Delta t}\) scaling is what distinguishes diffusion noise from ordinary deterministic increments.

\paragraph{Core intuition.}
Drift controls average evidence direction. Positive \(\mu\) pushes the process toward the upper boundary. Noise creates trial-to-trial variability. The boundary \(a\) controls the speed-accuracy tradeoff: a larger boundary requires more evidence, usually increasing accuracy but also increasing reaction time.

The DDM connects to Bayes through log-likelihood ratios. If observations arrive sequentially and each observation gives evidence for one of two hypotheses, the accumulated log-likelihood ratio is a decision variable. Crossing a threshold corresponds to sufficient posterior odds for one choice.

\paragraph{Main mechanism: log-likelihood ratio accumulation.}
Let hypotheses be \(H_+\) and \(H_-\). Observations \(Y_1,Y_2,\ldots\) are conditionally independent given the hypothesis. The posterior log odds after \(n\) observations are
\[
\log\frac{P(H_+\mid Y_{1:n})}{P(H_-\mid Y_{1:n})}.
\]
Using Bayes' theorem and conditional independence,
\[
\log\frac{P(H_+\mid Y_{1:n})}{P(H_-\mid Y_{1:n})}
=
\log\frac{P(H_+)}{P(H_-)}
+
\sum_{k=1}^n
\log\frac{p(Y_k\mid H_+)}{p(Y_k\mid H_-)}.
\]
Each observation adds one log-likelihood-ratio increment. When many small independent increments are approximated by a diffusion limit, the accumulated log odds become drift plus noise, which is the DDM structure.

\paragraph{Worked examples.}
\emph{Small numerical example.}
Suppose the prior odds are \(1\), so the initial log odds are \(0\). Three observations have likelihood ratios \(2,1.5,0.5\) in favor of \(H_+\). The accumulated log odds are
\[
\log 2+\log 1.5+\log 0.5
=\log(1.5)\approx0.405.
\]
The posterior odds are \(e^{0.405}=1.5\), so
\[
P(H_+\mid Y_{1:3})=\frac{1.5}{1+1.5}=0.6.
\]

\emph{Computational neuroscience example.}
In random-dot motion tasks, neural activity in parietal and frontal circuits has been modeled as accumulating motion evidence toward a choice threshold. The model links stimulus strength to drift, neural variability to diffusion, and boundary height to speed-accuracy tradeoff.

\emph{AI example.}
Sequential classifiers, early-exit neural networks, and active sensing systems use related logic: keep processing until confidence is high enough, then stop. The mathematical objects may be neural-network logits rather than spikes, but the log-odds accumulation idea remains useful.

\paragraph{Common misconceptions and failure modes.}
The DDM is not every decision process. It is a model for two-choice accumulation with specific assumptions. Boundary crossing is an abstraction, not necessarily a literal neuron hitting a hard threshold. Fitted drift, noise, and boundary parameters can trade off statistically, so model comparison and experimental design matter. The SDE notation \(dW_t\) is not an ordinary derivative; it represents Brownian increments with variance \(dt\).

\paragraph{Exercises.}
\begin{enumerate}[label=\textbf{50.1.\arabic*.}, leftmargin=*]
\item \textbf{Mechanical.} In a DDM with boundaries \(\pm a\), what choice is made if the path first hits \(-a\)? What parameter controls the average direction of evidence?
\item \textbf{Proof or derivation.} Derive the posterior log-odds sum from Bayes' theorem and conditional independence.
\item \textbf{Numerical/computational.} Simulate by hand three Euler steps with \(X_0=0\), \(\mu=1\), \(\sigma=0.5\), \(\Delta t=1\), and noise samples \(\xi=(0,-1,2)\).
\item \textbf{Interpretation.} Explain how increasing boundary height can increase accuracy and reaction time in the same model.
\item \textbf{Debug the misconception.} A student says, "The DDM proves decisions are Brownian motion in the brain." Correct the statement by distinguishing model abstraction from biological implementation.
\end{enumerate}

\subsection{50.2 Bayesian models of perception}

\paragraph{Motivating question.}
How can prior knowledge and noisy sensory data be combined into a posterior belief about a hidden cause?

\paragraph{Beginner setup.}
Perception often has an inverse structure. The world has hidden causes: object position, sound source, body state, illumination, motion direction. Sensory measurements are noisy effects of those causes. A generative model describes the forward direction:
\[
\text{latent cause}\to\text{observation}.
\]
Inference reverses the arrow:
\[
\text{observation}\to\text{belief about latent cause}.
\]

Bayes' theorem performs the reversal by reweighting and normalization:
\[
p(z\mid y)=\frac{p(y\mid z)p(z)}{p(y)}.
\]
The prior \(p(z)\) says which causes were plausible before the measurement. The likelihood \(p(y\mid z)\) says how likely the measurement is under each cause. The posterior \(p(z\mid y)\) is the updated belief.

\paragraph{Formal setup.}
Let \(Z\in\mathcal Z\) be a latent cause and \(Y\in\mathcal Y\) an observation. A generative model specifies
\[
p(z)\quad\text{and}\quad p(y\mid z).
\]
The joint distribution factors as
\[
p(z,y)=p(y\mid z)p(z).
\]
It also factors as
\[
p(z,y)=p(z\mid y)p(y).
\]
Equating these two factorizations gives Bayes' theorem:
\[
\boxed{
p(z\mid y)=\frac{p(y\mid z)p(z)}{p(y)}.
}
\]
For discrete \(z\),
\[
p(y)=\sum_{z'}p(y\mid z')p(z').
\]
For continuous \(z\),
\[
p(y)=\int p(y\mid z')p(z')\,dz'.
\]

Sequential Bayes uses the posterior after one observation as the prior for the next:
\[
\boxed{
p(z\mid y_{1:t})
\propto p(y_t\mid z)\,p(z\mid y_{1:t-1}).
}
\]
This is the prediction-update loop behind filtering.

\paragraph{Core intuition.}
Bayes is not magic confidence. It is reweighting. Hypotheses with high prior probability and high likelihood get larger posterior weight. The evidence normalizes all weights to sum or integrate to \(1\).

In perception, priors can be useful because sensory data are ambiguous. A dim object might be dark, far away, or poorly illuminated. A Bayesian model combines likelihood information from sensory measurements with prior structure learned from the environment or task.

\paragraph{Main mechanism: Gaussian cue combination.}
Suppose a scalar latent variable \(Z\) has prior
\[
Z\sim\mathcal N(\mu_0,\sigma_0^2),
\]
and a sensory observation satisfies
\[
Y\mid Z=z\sim\mathcal N(z,\sigma_y^2).
\]
After observing \(Y=y\), the posterior is Gaussian:
\[
Z\mid Y=y\sim\mathcal N(\mu_{\mathrm{post}},\sigma_{\mathrm{post}}^2),
\]
with precision
\[
\boxed{
\frac{1}{\sigma_{\mathrm{post}}^2}
=\frac{1}{\sigma_0^2}+\frac{1}{\sigma_y^2}
}
\]
and mean
\[
\boxed{
\mu_{\mathrm{post}}
=\sigma_{\mathrm{post}}^2
\left(
\frac{\mu_0}{\sigma_0^2}
+
\frac{y}{\sigma_y^2}
\right).
}
\]
The posterior mean is a precision-weighted average of prior mean and observation. More reliable cues have smaller variance, larger precision, and more influence.

\paragraph{Worked examples.}
\emph{Small numerical example.}
Let \(\mu_0=0\), \(\sigma_0^2=4\), \(y=3\), and \(\sigma_y^2=1\). The prior precision is \(1/4\) and the sensory precision is \(1\). Thus
\[
\frac{1}{\sigma_{\mathrm{post}}^2}=1.25,
\qquad
\sigma_{\mathrm{post}}^2=0.8.
\]
The posterior mean is
\[
\mu_{\mathrm{post}}=0.8\left(0/4+3/1\right)=2.4.
\]
The estimate moves toward the observation but remains pulled by the prior.

\emph{Computational neuroscience example.}
In multisensory perception, visual and vestibular cues to heading can be combined approximately by reliability weighting. A cue with lower variance should receive more weight. Neurally, this may correspond to populations whose response variability and tuning determine likelihood width, but the exact implementation is an empirical question.

\emph{AI example.}
Sensor fusion in robotics uses the same Gaussian update to combine a prior state estimate with a noisy measurement. Kalman filters, Bayesian neural networks, and probabilistic latent-variable models all use the principle that generative models produce observations and inference reverses that direction.

\paragraph{Common misconceptions and failure modes.}
A Bayesian model is only as good as its prior, likelihood, and assumptions. A prior is not automatically a bias in the bad sense; it is part of the model and can improve inference under uncertainty. However, a wrong prior can produce systematic errors. Bayesian perception models describe rational inference under specified assumptions; they do not automatically prove that human or animal perception is optimal.

\paragraph{Exercises.}
\begin{enumerate}[label=\textbf{50.2.\arabic*.}, leftmargin=*]
\item \textbf{Mechanical.} In Bayes' theorem \(p(z\mid y)=p(y\mid z)p(z)/p(y)\), identify the prior, likelihood, posterior, and evidence.
\item \textbf{Proof or derivation.} Derive Bayes' theorem from the two factorizations \(p(z,y)=p(y\mid z)p(z)=p(z\mid y)p(y)\).
\item \textbf{Numerical/computational.} Compute the Gaussian posterior mean and variance for \(\mu_0=1\), \(\sigma_0^2=9\), \(y=4\), and \(\sigma_y^2=1\).
\item \textbf{Interpretation.} Why does the lower-variance cue receive more weight in Gaussian cue combination?
\item \textbf{Debug the misconception.} A student says, "Bayes says the posterior equals the likelihood." Correct the statement using prior and evidence.
\end{enumerate}

\subsection{50.3 Predictive coding}

\paragraph{Motivating question.}
How can inference in a generative model be written as reducing precision-weighted prediction errors?

\paragraph{Beginner setup.}
Predictive coding begins with a generative model. Higher-level latent variables predict lower-level observations. The mismatch between prediction and observation is a prediction error. Inference updates latent estimates to reduce those errors, weighted by their reliability or precision.

The key point is modest and mathematical: for some probabilistic models, posterior inference is equivalent to minimizing an objective built from prediction-error terms. The stronger biological claim, that cortical circuits implement these updates in a particular way, requires additional evidence.

\paragraph{Formal setup.}
Consider a linear Gaussian generative model:
\[
Z\sim\mathcal N(\mu_z,\Sigma_z),
\qquad
Y\mid Z=z\sim\mathcal N(Cz,\Sigma_y).
\]
Let \(\Pi_z=\Sigma_z^{-1}\) and \(\Pi_y=\Sigma_y^{-1}\) be precision matrices. For an observed \(y\), define prediction errors
\[
\varepsilon_y=y-Cz,
\qquad
\varepsilon_z=z-\mu_z.
\]
The negative log posterior, up to constants independent of \(z\), is
\[
\boxed{
\mathcal F(z)
=\frac12\varepsilon_y^\top\Pi_y\varepsilon_y
+
\frac12\varepsilon_z^\top\Pi_z\varepsilon_z.
}
\]
Minimizing \(\mathcal F(z)\) gives the maximum a posteriori estimate of \(z\).

\paragraph{Core intuition.}
The observation \(y\) should be close to the prediction \(Cz\), but not all errors matter equally. If observation noise is small, \(\Pi_y\) is large and sensory prediction errors are weighted strongly. If the prior is reliable, \(\Pi_z\) is large and deviations from the prior mean are weighted strongly.

Predictive coding uses this precision-weighted error structure to describe iterative inference. Top-down predictions, bottom-up prediction errors, and precision weights are model variables. The biological mapping to cortical layers, cell types, or rhythms is a hypothesis, not a mathematical identity.

\paragraph{Main mechanism: gradient update for latent estimates.}
Differentiate
\[
\mathcal F(z)
=\frac12(y-Cz)^\top\Pi_y(y-Cz)
+
\frac12(z-\mu_z)^\top\Pi_z(z-\mu_z).
\]
The gradient is
\[
\nabla_z\mathcal F
=-C^\top\Pi_y(y-Cz)+\Pi_z(z-\mu_z).
\]
A gradient-descent inference update is
\[
\boxed{
\dot z
=-\nabla_z\mathcal F
=C^\top\Pi_y(y-Cz)-\Pi_z(z-\mu_z).
}
\]
Thus the latent estimate is pulled by sensory prediction error \(y-Cz\) and by the prior error \(z-\mu_z\), with both pulls weighted by precision.

\paragraph{Worked examples.}
\emph{Small numerical example.}
Let \(z,y\) be scalars, \(C=1\), \(\mu_z=0\), \(\Sigma_z=4\), \(\Sigma_y=1\), and \(y=3\). Then
\[
\mathcal F(z)=\frac12(3-z)^2+\frac12\frac{z^2}{4}.
\]
Setting the derivative to zero gives
\[
-(3-z)+\frac{z}{4}=0,
\]
so
\[
1.25z=3,\qquad z=2.4.
\]
This matches the Gaussian Bayes posterior mean from Section 50.2.

\emph{Computational neuroscience example.}
A hierarchical sensory model might posit that a higher area encodes a latent cause \(z\), predicts lower-level activity \(Cz\), and receives prediction errors when the lower-level activity differs. The equation above says how a latent estimate would update in a linear Gaussian model. Whether a specific cortical circuit implements these variables is a separate empirical claim.

\emph{AI example.}
Autoencoders, variational autoencoders, and self-supervised predictive models all compare predictions with observations. Predictive coding differs from ordinary prediction loss when it is framed as inference over latent causes under a generative model with precision-weighted errors. The mathematics overlaps with variational inference and optimization.

\paragraph{Common misconceptions and failure modes.}
Predictive coding is not just saying "the brain predicts." It requires a model of latent causes, predictions, errors, precision, and updates. It is not identical to backpropagation, although both use error signals and gradients. It should not be used to explain every neural response after the fact without specifying the generative model and falsifiable predictions. Precision weighting is not optional decoration; it determines which errors drive inference.

\paragraph{Exercises.}
\begin{enumerate}[label=\textbf{50.3.\arabic*.}, leftmargin=*]
\item \textbf{Mechanical.} For \(Y\mid Z=z\sim\mathcal N(Cz,\Sigma_y)\), write the sensory prediction error and its precision matrix.
\item \textbf{Proof or derivation.} Differentiate \(\mathcal F(z)=\frac12(y-Cz)^\top\Pi_y(y-Cz)+\frac12(z-\mu_z)^\top\Pi_z(z-\mu_z)\) to derive the predictive-coding update.
\item \textbf{Numerical/computational.} In the scalar example, recompute the MAP estimate when \(\Sigma_y=4\) and \(\Sigma_z=1\), with \(y=3\) and \(\mu_z=0\).
\item \textbf{Interpretation.} What happens to the update if sensory precision \(\Pi_y\) is very small?
\item \textbf{Debug the misconception.} A student says, "Predictive coding proves cortex performs exact Bayesian inference." Correct the statement by separating the mathematical model from biological implementation.
\end{enumerate}

\subsection{50.4 Active inference and control}

\paragraph{Motivating question.}
How do beliefs, predictions, preferences, and actions interact when an agent can change the observations it receives?

\paragraph{Beginner setup.}
Passive perception updates beliefs from observations. Active agents also act. An action changes the future state of the world or the future sensory sample. For example, moving the eyes changes visual input; reaching changes proprioceptive and tactile input; asking a question changes future information.

Active inference is one framework for modeling this closed loop. It uses a generative model over hidden states, observations, and actions or policies. The agent evaluates policies by predicted outcomes, preferences, uncertainty, and information gain. This is closely related to partially observed control and model-based reinforcement learning, but uses variational-inference language.

\paragraph{Formal setup.}
Let hidden states be \(X_t\), observations \(O_t\), actions \(A_t\), and policies
\[
\pi=(a_t,a_{t+1},\ldots,a_T).
\]
A simple controlled generative model factors as
\[
p(x_{1:T},o_{1:T}\mid \pi)
=p(x_1)\prod_{t=1}^T p(o_t\mid x_t)
\prod_{t=1}^{T-1}p(x_{t+1}\mid x_t,a_t).
\]
Inference estimates beliefs such as
\[
q(x_t\mid o_{1:t},\pi).
\]
Policy evaluation uses predicted future observations and states under \(\pi\).

A beginner-friendly decomposition of expected free energy is
\[
\boxed{
G(\pi)
\approx
\underbrace{\mathbb E_{q(o\mid\pi)}[-\log p_{\mathrm{pref}}(o)]}_{\text{expected preference cost}}
-
\underbrace{\mathbb E_{q(o\mid\pi)}
\left[
D_{\mathrm{KL}}(q(x\mid o,\pi)\|q(x\mid\pi))
\right]}_{\text{expected information gain}}.
}
\]
Different texts use related decompositions with risk, ambiguity, and epistemic value. The common idea is that policies are favored when they are expected to reach preferred observations and, in some settings, reduce uncertainty about hidden states.

\paragraph{Core intuition.}
Action can serve two roles. It can exploit: choose actions expected to produce preferred outcomes. It can explore: choose actions expected to reveal hidden state information. A good model-based agent often needs both.

The active-inference vocabulary differs from standard optimal control, but the mathematical ingredients are familiar: hidden states, observations, transition models, preferences or costs, policies, and approximate inference. The careful question is what objective is being optimized and what approximations are used.

\paragraph{Main mechanism: a two-action information-seeking example.}
Suppose an agent must choose between two actions before making a final decision. The hidden state \(X\in\{0,1\}\) is equally likely. Action \(a_{\mathrm{look}}\) gives an observation that reveals \(X\) correctly with probability \(0.9\) but has no immediate reward. Action \(a_{\mathrm{guess}}\) gives no information but allows an immediate guess with success probability \(0.5\).

The information gain from a future observation \(O\) can be measured by
\[
I(X;O\mid a)
=H(X)-H(X\mid O,a).
\]
For \(a_{\mathrm{guess}}\), \(O\) carries no information, so \(I(X;O\mid a_{\mathrm{guess}})=0\). For \(a_{\mathrm{look}}\), \(O\) is informative, so \(H(X\mid O,a_{\mathrm{look}})<H(X)\) and the information gain is positive. If future reward depends strongly on knowing \(X\), looking may be preferred even without immediate reward.

This illustrates epistemic action: action chosen partly to reduce uncertainty. It does not require claiming that all action is information seeking; preferences, costs, and task demands also matter.

\paragraph{Worked examples.}
\emph{Small numerical example.}
With \(P(X=0)=P(X=1)=1/2\), \(H(X)=1\) bit. If \(a_{\mathrm{look}}\) produces a binary observation that is correct with probability \(0.9\), then after either observation the posterior is \((0.9,0.1)\). The conditional entropy is
\[
H(X\mid O,a_{\mathrm{look}})
=-0.9\log_2(0.9)-0.1\log_2(0.1)
\approx0.469.
\]
So
\[
I(X;O\mid a_{\mathrm{look}})
\approx1-0.469=0.531\ \text{bits}.
\]

\emph{Computational neuroscience example.}
Eye movements can be modeled as actions that sample informative parts of a visual scene while also serving task goals. An active-inference model would specify hidden scene states, observations generated by fixation, action-dependent transitions, and preferences over outcomes. The model should then be compared with behavior and neural data.

\emph{AI example.}
Partially observable reinforcement learning, active sensing, Bayesian experimental design, and model-based exploration all contain similar mathematics. Policies are evaluated not only by immediate reward but also by how they change future belief states. Active inference is one formulation of this broad model-based control problem.

\paragraph{Common misconceptions and failure modes.}
Active inference is not a license to say that every action minimizes a vague free energy. The hidden states, observations, policies, preferences, and approximate posterior must be specified. Expected free energy is not identical to thermodynamic free energy, although the names are historically connected. Biological applications should distinguish descriptive curve fitting from mechanistic evidence. In engineering, standard control or reinforcement learning may be simpler and more transparent for many tasks.

\paragraph{Exercises.}
\begin{enumerate}[label=\textbf{50.4.\arabic*.}, leftmargin=*]
\item \textbf{Mechanical.} In the controlled generative model, identify which factor describes observations and which factor describes action-dependent state transitions.
\item \textbf{Proof or derivation.} Show that if an action's observation is independent of \(X\), then \(I(X;O\mid a)=0\).
\item \textbf{Numerical/computational.} Compute the information gain when a binary observation is correct with probability \(0.8\) instead of \(0.9\), with a uniform binary hidden state.
\item \textbf{Interpretation.} Give one example of an action that is valuable mainly because it gathers information.
\item \textbf{Debug the misconception.} A student says, "Active inference proves organisms act only to reduce uncertainty." Correct the statement using preferences, costs, and model assumptions.
\end{enumerate}

\paragraph{Closing bridge.}
Decision and perception models combine dynamics, probability, and optimization. Drift-diffusion models accumulate evidence; Bayesian perception reweights latent causes; predictive coding turns Gaussian inference into precision-weighted error minimization; active inference adds policy-dependent observations and preferences. Chapter 51 will treat such models as objects to fit, compare, criticize, and validate against neural and behavioral data.

\clearpage
